\chapter*{导\quad 论}
\phantomsection\addcontentsline{toc}{chapter}{导\quad 论}
\markboth{导\quad 论}{导\quad 论}

一九六三年春，莫斯科，苏联科学院远东委员会设于一栋旧楼三层的档案室，一名读汉文的研究生被叫去看三只木箱。编目卡上的类别写作「远东出土文书」，箱内是三十七片长约二寸至九寸、宽约一指的薄片，多为金属，间有陶瓷，以棉纸逐片隔开，片上有刻字，字为小篆。此后的六十年里，这批材料经过三次编目、一次转移、两次被列入清退名单，始终没有人调阅。使它搁置的原因与保密无关，只是没有人知道应当把它归入哪一类：从材质看是器物，从字迹看是文书，从出土地点看则不属于任何一个国家的考古区划。一九九二年基金会接管远东委员会遗留卷宗之后，它被编入 \textsc{фв} 系列，编号自 \textsc{фв}-1963-11 起，与一批来历相近、彼此并无关联的散件混在一处。本书所处理的，正是这三十七片。

\section*{一、材料的性质}

关于这批材料的出处，档案仅有一行著录：
\begin{center}
\small\textsc{фв}-1963-11\quad 月球正面，\textsc{mare-tranq-07} 基地，第 4 号居住环，西侧储物间
\end{center}
\textsc{mare-tranq-07} 即静海七号（见第 117 全宗第 4 目第 2187 号卷宗第 12 叶及《月面设施复原图集》图版 11），亦即汉译「广寒宫」所指之地。同批材料内另存两张照片，摄于取走残片的那个房间，画面所示西壁架数与简七的自记相合，可证三十七片确出自该处。

这批材料在史料学上的位置需要先行界定。它属于原件，并非从他书引文中钩沉出来的佚文：字口清晰，可辨部分约占十分之六七，残损处蚀去的是字，不是载体。它同样不属于传统意义上的出土文书，因为出土地点不在中国境内，而在月球正面的月海之中。这两点使其与常规的辑佚工作区别开来。苏联方面在一九六三年曾对同类材料作过一次成分分析，报告编号 \textsc{фв}-1963-07，结论部分只留下一句「来源不明，暂不归类」。这句话写在编目卡上，不在原件上，却准确地概括了此后六十年间的处境。

\section*{二、材料的内容}

三十七片构不成一部完整的书，而是一处设施在长时段内陆续刻写的簿籍与规程：居住环的形制、气路与水路的分布、人口的编制、口粮的定额、名册的更替、事故的数目，以及撤离前的清点与退役。其中分量最重的是卷三的十二片，所记为公元前二百二十一年前后的一次清点，与莫斯科档案馆第 117 全宗第 2187 号卷宗第 121 叶所录「设施退役清单（共十九项）」同出一源；俄文档案抄本仅摘引了第七项，本书给出全文。就内容而言，这批材料所记的是一处基地的日常行政，而非任何宏大事件：三十七片记下的是岁一检的管道、岁一焚的名册、日三合的定额，以及最后一年里逐条执行的十九项处置；人名、战争、决策过程均无记载，唯一一处带有动作的记述，仅有简三十一「闭者七人，回首一次」。

就材料所及，可以为此城勾勒出一个基本轮廓。它在月球正面静海西缘，直径约一点二公里，六环相套，居民十一万九千人。居民不属于地球人类群类，其体质与地上人异趣：不老化，不生育，除外伤与极大量的污秽之外不会死亡。城的洁净不靠天然，而依赖一套藏于外墙之内的系统常年运转，系统的能源取自地下。公元前二百二十一年，地下热源降至临界以下，全城分六批迁离，城池随即封存。以上数语皆为下编各章的结论，此处提前写出，只为使读者在进入释文之前明晰这批文字所记为何。各条的依据与限定分见下编，尤其是第二章、第六章与第七章；凡传世文献与此城相关之事，本书一概不涉，其故见下编前言。行政记录不修辞，不表彰，不隐恶，这一特征既是材料的局限，也是它的效用所在：它只记下当时的人认为需要留下来的东西，而何物被认作需要留下来，本身即构成一项可资研究的对象。

\section*{三、体例}

本书分上下两编。上编为释文，按内容分作营建、岁计、撤离、后题四卷，另附存疑八条，下编一概不予置论。每片分三段：释文照原件行款录写，出处记编号、尺寸与材质，校记则专司字词之考订；校记不作论断，凡涉及释读之判断，一律移入下编。下编为考释，共八章：第一章论月之民的体质与制度，第二章据简六所记尺度复原城的形制，第三章论净秽制度，第四章仅处理简二十二第十三项「药，弃」二字，第五章论净秽山与月兔，第六章处理岁计与人口之数目，第七章还原撤离的经过，第八章论名与籍。

两编的读法有所不同：上编可以跳读，按下编索引所举简次检阅即可；下编则须顺序通读，因其结论系逐层推演而成，凡例另列于上编之前，不复赘述。尚有一事须于此处交代：原件所刻为小篆，释文一律改作今字，此乃整理工作之通例，并非原件之字形；书中插图亦照释文摹写。若欲见原字形，须待影印本，影印本拟随第三版附出，此次未及收入。

\section*{四、方法与限度}

本书的编次出自笔者。原件的原始编号与摆放次序均已无从确认，旧照上尚可辨识者仅剩层数；现行编次依据材质、字口与内容三类线索的交叉印证，其中任何一类单独成立都不足以定序，换作他人整理，次序亦将不同。释文的情况与此相类：在可辨的十分之六七之中，约有一成为笔者所补，凡补字一概用方框标出，并在校记中逐条说明依据；原件尚在，读者尽可持异议或重行补定。

数字方面，凡能互相校验者均已核校，其无法弥合者径直保留于正文之内，不加弥缝。其中最显著者为八千之差额：完工时十五万人（见第 117 全宗第 2187 号卷宗及俄方先期复原），六批共迁出十二万人，两万年间减员三万；简十九记事故二万二千，两数相减仍余八千，此八千人在残简中无一字记载。本书第六章列出三种可能解释，笔者倾向前第一种，并明言缺乏直接证据。至于年表问题，城中自有纪年，与中原不同而无法换算；残简中提及城中纪年者三处（十四年、十九年、二十三年），另有一处作「元年」，前者与后者是否同属一纪年无从判定，故附录二将其分列。

\section*{五、自白}

笔者不认为本书的编次是唯一正确的方案，只认为它是一种必要的编次。三十七片放在一间屋子里六十年，没有人调阅，将其整理成书并不指望扭转这一境遇；这批材料此后多半仍将滞留于原库房内，至多不过更换一个编目号。若谓此项工作有何效用，无非是将这三十七片所刻的文字从金属与陶瓷移至纸面；纸张不褪色，亦不腐烂，只要不遭焚毁，这些记录便得以长久存留于彼处，待后人翻检。

\begin{flushright}
{\small 巴兰尼科夫}\quad{\footnotesize\color{pencil} 二〇二四年六月，于莫斯科}
\end{flushright}
